% Options for packages loaded elsewhere
% Options for packages loaded elsewhere
\PassOptionsToPackage{unicode}{hyperref}
\PassOptionsToPackage{hyphens}{url}
\PassOptionsToPackage{dvipsnames,svgnames,x11names}{xcolor}
%
\documentclass[
  11pt,
  a4paper]{article}
\usepackage{xcolor}
\usepackage[left=1in,top=1in,right=1in,bottom=1in]{geometry}
\usepackage{amsmath,amssymb}
\setcounter{secnumdepth}{-\maxdimen} % remove section numbering
\usepackage{iftex}
\ifPDFTeX
  \usepackage[T1]{fontenc}
  \usepackage[utf8]{inputenc}
  \usepackage{textcomp} % provide euro and other symbols
\else % if luatex or xetex
  \usepackage{unicode-math} % this also loads fontspec
  \defaultfontfeatures{Scale=MatchLowercase}
  \defaultfontfeatures[\rmfamily]{Ligatures=TeX,Scale=1}
\fi
\usepackage{lmodern}
\ifPDFTeX\else
  % xetex/luatex font selection
\fi
% Use upquote if available, for straight quotes in verbatim environments
\IfFileExists{upquote.sty}{\usepackage{upquote}}{}
\IfFileExists{microtype.sty}{% use microtype if available
  \usepackage[]{microtype}
  \UseMicrotypeSet[protrusion]{basicmath} % disable protrusion for tt fonts
}{}
\makeatletter
\@ifundefined{KOMAClassName}{% if non-KOMA class
  \IfFileExists{parskip.sty}{%
    \usepackage{parskip}
  }{% else
    \setlength{\parindent}{0pt}
    \setlength{\parskip}{6pt plus 2pt minus 1pt}}
}{% if KOMA class
  \KOMAoptions{parskip=half}}
\makeatother
% Make \paragraph and \subparagraph free-standing
\makeatletter
\ifx\paragraph\undefined\else
  \let\oldparagraph\paragraph
  \renewcommand{\paragraph}{
    \@ifstar
      \xxxParagraphStar
      \xxxParagraphNoStar
  }
  \newcommand{\xxxParagraphStar}[1]{\oldparagraph*{#1}\mbox{}}
  \newcommand{\xxxParagraphNoStar}[1]{\oldparagraph{#1}\mbox{}}
\fi
\ifx\subparagraph\undefined\else
  \let\oldsubparagraph\subparagraph
  \renewcommand{\subparagraph}{
    \@ifstar
      \xxxSubParagraphStar
      \xxxSubParagraphNoStar
  }
  \newcommand{\xxxSubParagraphStar}[1]{\oldsubparagraph*{#1}\mbox{}}
  \newcommand{\xxxSubParagraphNoStar}[1]{\oldsubparagraph{#1}\mbox{}}
\fi
\makeatother


\usepackage{longtable,booktabs,array}
\newcounter{none} % for unnumbered tables
\usepackage{calc} % for calculating minipage widths
% Correct order of tables after \paragraph or \subparagraph
\usepackage{etoolbox}
\makeatletter
\patchcmd\longtable{\par}{\if@noskipsec\mbox{}\fi\par}{}{}
\makeatother
% Allow footnotes in longtable head/foot
\IfFileExists{footnotehyper.sty}{\usepackage{footnotehyper}}{\usepackage{footnote}}
\makesavenoteenv{longtable}
\usepackage{graphicx}
\makeatletter
\newsavebox\pandoc@box
\newcommand*\pandocbounded[1]{% scales image to fit in text height/width
  \sbox\pandoc@box{#1}%
  \Gscale@div\@tempa{\textheight}{\dimexpr\ht\pandoc@box+\dp\pandoc@box\relax}%
  \Gscale@div\@tempb{\linewidth}{\wd\pandoc@box}%
  \ifdim\@tempb\p@<\@tempa\p@\let\@tempa\@tempb\fi% select the smaller of both
  \ifdim\@tempa\p@<\p@\scalebox{\@tempa}{\usebox\pandoc@box}%
  \else\usebox{\pandoc@box}%
  \fi%
}
% Set default figure placement to htbp
\def\fps@figure{htbp}
\makeatother


% definitions for citeproc citations
\NewDocumentCommand\citeproctext{}{}
\NewDocumentCommand\citeproc{mm}{%
  \begingroup\def\citeproctext{#2}\cite{#1}\endgroup}
\makeatletter
 % allow citations to break across lines
 \let\@cite@ofmt\@firstofone
 % avoid brackets around text for \cite:
 \def\@biblabel#1{}
 \def\@cite#1#2{{#1\if@tempswa , #2\fi}}
\makeatother
\newlength{\cslhangindent}
\setlength{\cslhangindent}{1.5em}
\newlength{\csllabelwidth}
\setlength{\csllabelwidth}{3em}
\newenvironment{CSLReferences}[2] % #1 hanging-indent, #2 entry-spacing
 {\begin{list}{}{%
  \setlength{\itemindent}{0pt}
  \setlength{\leftmargin}{0pt}
  \setlength{\parsep}{0pt}
  % turn on hanging indent if param 1 is 1
  \ifodd #1
   \setlength{\leftmargin}{\cslhangindent}
   \setlength{\itemindent}{-1\cslhangindent}
  \fi
  % set entry spacing
  \setlength{\itemsep}{#2\baselineskip}}}
 {\end{list}}
\usepackage{calc}
\newcommand{\CSLBlock}[1]{\hfill\break\parbox[t]{\linewidth}{\strut\ignorespaces#1\strut}}
\newcommand{\CSLLeftMargin}[1]{\parbox[t]{\csllabelwidth}{\strut#1\strut}}
\newcommand{\CSLRightInline}[1]{\parbox[t]{\linewidth - \csllabelwidth}{\strut#1\strut}}
\newcommand{\CSLIndent}[1]{\hspace{\cslhangindent}#1}



\setlength{\emergencystretch}{3em} % prevent overfull lines

\providecommand{\tightlist}{%
  \setlength{\itemsep}{0pt}\setlength{\parskip}{0pt}}



 


\usepackage{iftex}

% % Minion Pro
\usepackage[lf]{MinionPro}
\usepackage{MyriadPro}

% IBM Plex
% \usepackage{fontspec}
% \usepackage{unicode-math}
% \setmainfont{IBM Plex Serif}
% \setsansfont{IBM Plex Sans}
% \setmonofont{IBM Plex Mono}
% \setmathfont{IBM Plex Math Beta 240212}

% STIX Two
% \usepackage{fontspec}
% \setmainfont{Stix Two Text}
% \setmathfont{Stix Two Math}
% \setsansfont[Scale=0.9608140457]{Fira Sans} % 7.2489/7.54454

% Libertinus Serif + Libertinus Math + Source Sans Pro + Source Code Pro
% \ifPDFTeX
% \usepackage[mono=false,osf]{libertine}
% \usepackage[libertine]{newtxmath}
% % \usepackage[scale=1,osf]{sourcesanspro} % 7.1832/7.1832
% % \usepackage[scale=0.8772671637]{sourcecodepro} % 7.1832/7.36934*0.9
% \usepackage[scale=0.9181263229,osf]{roboto} % 7.1832/7.82376
% \usepackage[scale=0.8132241635]{roboto-mono} % 7.1832/7.94969*0.9
% \else
% \usepackage[mono=false,osf]{libertine}
% \setmathfont[Scale=MatchUppercase]{libertinusmath-regular.otf}
% % \usepackage[scale=0.9850312868,osf]{sourcesanspro} % 7.2051/7.31459
% % \usepackage[scale=0.8865281581]{sourcecodepro} % 7.2051/7.31459*0.9
% \usepackage[scale=0.9126225152,osf]{roboto} % 7.2051/7.89494
% \usepackage[scale=0.8213602637]{roboto-mono} % 7.2051/7.89494*0.9
% \fi

% Enable section and paragraph numbering
\setcounter{secnumdepth}{4} % Numbering down to paragraph level
\setcounter{tocdepth}{4} % Include paragraphs in the ToC

%%% CUSTOMIZE SECTION SIZES
\makeatletter
\renewcommand{\section}{%
  \@startsection{section}{1}
    {0pt} 					% indent
    {1.0em plus 0.2em minus 0.2em} 	% beforeskip
    {0.001em plus 0.001em minus 0.001em}    	% afterskip
    {\Large\bfseries\sffamily}%
}
\renewcommand{\subsection}{%
  \@startsection{subsection}{2}
    {0pt} 					% indent
    {1.0em plus 0.2em minus 0.2em} 	% beforeskip
    {0.001em plus 0.001em minus 0.001em}    	% afterskip
    {\large\bfseries\sffamily}%
}
\renewcommand{\subsubsection}{%
  \@startsection{subsubsection}{3}
    {0pt} 					% indent
    {1.0em plus 0.2em minus 0.2em} 	% beforeskip
    {0.001em plus 0.001em minus 0.001em}    	% afterskip
    {\normalsize\bfseries\sffamily}%
}
\renewcommand{\paragraph}{%
  \@startsection{paragraph}{4}
    {0pt} 					% indent
    {1.0em plus 0.2em minus 0.2em} 	% beforeskip
    {0.001em plus 0.001em minus 0.001em}    	% afterskip
    {\normalsize\bfseries\sffamily}%
}
\makeatother

%%% GENERAL PACKAGES
\usepackage{hyphenat} % don't hyphenat titles
\usepackage{graphicx}
\usepackage{tabularx} 
\usepackage{marvosym} % Mail logo
\usepackage{enumitem} % Enumerate items (list)
\usepackage{amsmath} % for enhanced math features
\usepackage[labelfont=bf,justification=justified,singlelinecheck=false]{caption} % Figures caption setup
\usepackage[rightcaption]{sidecap} % figure with sidecaption
\usepackage[]{lineno} % Line numbers
\usepackage{ragged2e} % justify
\usepackage{array}
\usepackage{lscape} % landscape tables
\usepackage{rotating} % landscape tables
\usepackage{bm} % bold math

%%%% ORCID LOGO
\usepackage{academicons}
\definecolor{orcidlogocol}{HTML}{A6CE39}
\usepackage{orcidlink}

%%%% LAYOUT 
\tolerance=200
\emergencystretch=2em
\hyphenpenalty=1000
\hbadness=2000
\widowpenalty=1000 % widow penalty
\clubpenalty=1000 % orphans penalty
\renewcommand{\baselinestretch}{1.3333333333333333333333333333} % line spacing

\renewcommand{\arraystretch}{1.3333}
\makeatletter
\@ifpackageloaded{float}{}{\usepackage{float}}
\floatstyle{plain}
\@ifundefined{c@chapter}{\newfloat{suppfig}{h}{losuppfig}}{\newfloat{suppfig}{h}{losuppfig}[chapter]}
\floatname{suppfig}{Figure S}
\newcommand*\quartosuppfigref[1]{Figure \hyperref[#1]{S\ref{#1}}}
\@ifpackageloaded{caption}{}{\usepackage{caption}}
\DeclareCaptionLabelFormat{quartosuppfigreflabelformat}{#1#2}
\captionsetup[suppfig]{labelformat=quartosuppfigreflabelformat}
\newcommand*\listofsuppfigs{\listof{suppfig}{List of Supplementary Figures}}
\makeatother
\makeatletter
\@ifpackageloaded{float}{}{\usepackage{float}}
\floatstyle{plain}
\@ifundefined{c@chapter}{\newfloat{supptbl}{h}{losupptbl}}{\newfloat{supptbl}{h}{losupptbl}[chapter]}
\floatname{supptbl}{Table S}
\floatstyle{plaintop}
\restylefloat{supptbl}
\newcommand*\quartosupptblref[1]{Table \hyperref[#1]{S\ref{#1}}}
\@ifpackageloaded{caption}{}{\usepackage{caption}}
\DeclareCaptionLabelFormat{quartosupptblreflabelformat}{#1#2}
\captionsetup[supptbl]{labelformat=quartosupptblreflabelformat}
\newcommand*\listofsupptbls{\listof{supptbl}{List of Supplementary Tables}}
\makeatother
\makeatletter
\@ifpackageloaded{caption}{}{\usepackage{caption}}
\AtBeginDocument{%
\ifdefined\contentsname
  \renewcommand*\contentsname{Table of contents}
\else
  \newcommand\contentsname{Table of contents}
\fi
\ifdefined\listfigurename
  \renewcommand*\listfigurename{List of Figures}
\else
  \newcommand\listfigurename{List of Figures}
\fi
\ifdefined\listtablename
  \renewcommand*\listtablename{List of Tables}
\else
  \newcommand\listtablename{List of Tables}
\fi
\ifdefined\figurename
  \renewcommand*\figurename{Figure}
\else
  \newcommand\figurename{Figure}
\fi
\ifdefined\tablename
  \renewcommand*\tablename{Table}
\else
  \newcommand\tablename{Table}
\fi
}
\@ifpackageloaded{float}{}{\usepackage{float}}
\floatstyle{ruled}
\@ifundefined{c@chapter}{\newfloat{codelisting}{h}{lop}}{\newfloat{codelisting}{h}{lop}[chapter]}
\floatname{codelisting}{Listing}
\newcommand*\listoflistings{\listof{codelisting}{List of Listings}}
\makeatother
\makeatletter
\makeatother
\makeatletter
\@ifpackageloaded{caption}{}{\usepackage{caption}}
\@ifpackageloaded{subcaption}{}{\usepackage{subcaption}}
\makeatother
\usepackage{bookmark}
\IfFileExists{xurl.sty}{\usepackage{xurl}}{} % add URL line breaks if available
\urlstyle{same}
\hypersetup{
  pdftitle={Mechanical power output during stretch--shortening cycles of rat medial gastrocnemius muscle: influence of various muscle length trajectories -- S2 File},
  colorlinks=true,
  linkcolor={blue},
  filecolor={Maroon},
  citecolor={Blue},
  urlcolor={Blue},
  pdfcreator={LaTeX via pandoc}}


\title{Mechanical power output during stretch--shortening cycles of rat
medial gastrocnemius muscle: influence of various muscle length
trajectories -- S2 File}
\author{}
\date{}
\begin{document}


\linenumbers

\renewcommand{\thesection}{S\arabic{section}}
\setcounter{section}{1}
\floatname{suppfig}{Figure }
\renewcommand{\thesuppfig}{\thesection.\arabic{suppfig}}
\renewcommand*\quartosuppfigref[1]{Figure~\ref{#1}}

\phantomsection

\section{Optimisation protocol to identify optimal MTC length over
time}\label{optimisation-protocol-to-identify-optimal-mtc-length-over-time}

One of the aims of this study was to identify the optimal MTC length
over time for maximal AMPO, when MTC length over time is entirely free
to take any form. To predict this, we used a Hill-type MTC model. In our
Hill-type MTC model, we assumed that both SEE and PEE were purely
elastic elements (e.g., \citeproc{ref-anderson_dynamic_1999}{Anderson
and Pandy, 1999}; \citeproc{ref-van_soest_contribution_1993}{Soest and
Bobbert, 1993}). Therefore, during fully periodic SSC neither PEE nor
SEE contributed to AMPO of the MTC. This also means that neither PEE nor
SEE influences the optimal CE length over time that yields maximal AMPO
over a full periodic SSC. We used this to simplify the optimisation
problem. We reformulated to model such that we considered a Hill-type
model of CE only. As such, CE length was no longer a state variable, but
an input to the model. The model therefore had two inputs: CE length
(\(L_{CE}\)) and muscle stimulation (\(STIM\)) over time. The only state
variable was \(\gamma\), the normalised intrafilament \(Ca^{2+}\)
concentration. Given this formulation, \(L_{CE}\), \(STIM\) and
\(\gamma\) were known at each time instant during the optimisation
process. Consequently, CE velocity can be calculated by taking the
time-derivative of CE length. Using CE velocity, CE force could then be
computed using the CE force--velocity relationship. MTC length over time
was subsequently obtained by taking the following steps: 1) PEE force
was calculated from PEE length (which equals CE length); 2) SEE force
was calculated by taking the sum of CE and PEE force; 3) SEE length was
calculated from SEE force; 4) MTC length was calculated as the sum of CE
and SEE length. By computing MTC length over time via these steps, we
could also impose MTC length excursion.

The optimisation problem was thus to find the periodic CE length,
\(STIM\) and \(\gamma\) over time for maximal AMPO. Optimisations were
performed for 1) imposed cycle frequencies (1.0--6.0 Hz in 0.5 Hz
steps); 2) imposed FTS values (0.05--0.95 in 0.05 steps) and 3) imposed
MTC length excursions (2--11\,mm in 1\,mm steps). For all optimisation,
the following cost function was maximised (note that, in our
formulation, CE shortening corresponds to negative CE velocity):

\begin{equation}\protect\phantomsection\label{eq-objective}{
J = -\int_{t_0}^{t_0+T_c} F_{CE}^{rel} \cdot v_{CE}^{rel} \ dt
}\end{equation} where \(t_0\) denotes the initial time and \(T_c\)
denotes the cycle time.

The solution had to satisfy the differential equation describing the
excitation dynamics:
\begin{equation}\protect\phantomsection\label{eq-dynamic-con}{
\dot{\gamma} = f(\gamma(t),STIM(t))
}\end{equation}

The following task constraints were added to achieve periodic behaviour:
\begin{equation}\protect\phantomsection\label{eq-task-con}{
\begin{gathered}
  L_{CE}(t_0) = L_{CE}(t_0+T_c) \\
  \gamma(t_0) = \gamma(t_0+T_c)
\end{gathered}
}\end{equation}

Moreover, inequality constraints were imposed on \(\gamma(t)\),
\(L_{CE}^{rel}(t)\) and \(STIM(t)\):
\begin{equation}\protect\phantomsection\label{eq-ineq-con}{
\begin{gathered}
  \gamma_0 \leq \gamma(t) \leq 1 \\
  (1-w) \leq L_{CE}^{rel}(t) \leq (1+w) \\
  0 \leq STIM(t) \leq 1
\end{gathered}
}\end{equation} where \(L_{CE}^{rel}\) denotes the CE length normalised
to CE optimum length (i.e.~the CE length yielding maximal isometric CE
force).

We implemented our problem as a Multiphase Optimal Control Problem to
assure that CE shortening and lengthening occurred only once per full
periodic cycle, which was done by inequality constraints on
\(v_{CE}^{rel}(t)\):
\begin{equation}\protect\phantomsection\label{eq-fts-con}{
\begin{aligned}
  v_{CE}^{rel}(t_0 \ ... \ t_0 + T_c \cdot FTS) &\leq 0 \\
  v_{CE}^{rel}(t_0 + T_c \cdot FTS \ ... \ t_0 + T_c) &\geq 0
\end{aligned}
}\end{equation} FTS was a numerical value (between 0.05 and 0.95) when
FTS were imposed (see above). When FTS was not imposed, FTS was a
control variable optimised for maximal AMPO.

When MTC length excursion was imposed, an extra inequality constraint
was added:\\
\begin{equation}\protect\phantomsection\label{eq-lcecon}{
  L_{MTC}(t_0) + L_{MTC}(t_0 + T_c \cdot FTS) - \textrm{MTC length excursion}
}\end{equation}

The dynamic, task and (in)equality constraints were implemented in
CasADi (\citeproc{ref-andersson_casadi_2019}{Andersson et al., 2019})
and the cost function was minimised using a Direct Collocation method
using a nonlinear interior point method (IPOPT,
\citeproc{ref-wachter_implementation_2006}{Wächter and Biegler, 2006}).
The optimal control solution was checked by running forward simulations
of the model with a linear-interpolated \(F_{CE}^{rel}(t)\) and
\(STIM(t)\) as inputs to confirm that no relevant constraints were
violated between the collocation points. This was done using SciPy's ODE
integrator with an Implicit Runge-Kutta method of Radau IIA family of
order 5 (\citeproc{ref-hairer_implementation_1996}{Hairer and Wanner,
1996}) and an absolute and relative tolerance of 1e-6 and a maximum
timestep of 1e-3. This yielded only very small differences between the
states at the collocation points and those at the same time instance of
the forward simulation for all optimisations.

\section*{References}\label{references}
\addcontentsline{toc}{section}{References}

\protect\phantomsection\label{refs}
\begin{CSLReferences}{1}{1}
\bibitem[\citeproctext]{ref-anderson_dynamic_1999}
\textbf{Anderson, F. C. and Pandy, M. G.} (1999).
\href{https://doi.org/10.1080/10255849908907988}{A {Dynamic}
{Optimization} {Solution} for {Vertical} {Jumping} in {Three}
{Dimensions}}. \emph{Computer Methods in Biomechanics and Biomedical
Engineering} \textbf{2}, 201--231.

\bibitem[\citeproctext]{ref-andersson_casadi_2019}
\textbf{Andersson, J. A. E., Gillis, J., Horn, G., Rawlings, J. B. and
Diehl, M.} (2019).
\href{https://doi.org/10.1007/s12532-018-0139-4}{{CasADi}: A software
framework for nonlinear optimization and optimal control}.
\emph{Mathematical Programming Computation} \textbf{11}, 1--36.

\bibitem[\citeproctext]{ref-hairer_implementation_1996}
\textbf{Hairer, E. and Wanner, G.} (1996).
\href{https://doi.org/10.1007/978-3-642-05221-7_8}{Implementation of
{Implicit} {Runge}-{Kutta} {Methods}}. In \emph{Solving {Ordinary}
{Differential} {Equations} {II}: {Stiff} and {Differential}-{Algebraic}
{Problems}} (ed. Hairer, E.) and Wanner, G.), pp. 118--130. Berlin,
Heidelberg: Springer Berlin Heidelberg.

\bibitem[\citeproctext]{ref-van_soest_contribution_1993}
\textbf{Soest, A. J. van and Bobbert, M. F.} (1993).
\href{https://doi.org/10.1007/BF00198959}{The contribution of muscle
properties in the control of explosive movements}. \emph{Biological
Cybernetics} \textbf{69}, 195--204.

\bibitem[\citeproctext]{ref-wachter_implementation_2006}
\textbf{Wächter, A. and Biegler, L. T.} (2006).
\href{https://doi.org/10.1007/s10107-004-0559-y}{On the implementation
of an interior-point filter line-search algorithm for large-scale
nonlinear programming}. \emph{Mathematical Programming} \textbf{106},
25--57.

\end{CSLReferences}




\end{document}
